\section{Accuracy and resource comparisons: derivations}
\label{app:comparison}

We compare approximations to the same population predictor $f_\infty$ of
\eqref{eq:dense-flow}, assumed to exist on $[0,T]$, with the same data and
initialization law, including the readout.  Matching initial predictions alone
does not match the initial tangent kernel.  Throughout, constants can depend
on $m,H,T$, the data, initialization and test law; displayed factors in these
variables do not exhaust that dependence.  Population error estimates below
are stated hypotheses, whereas the fixed-width response-memory results are
\Cref{thm:old,thm:new}.  We suppress the $T$ subscripts on the error constants
of \Cref{sec:resources}.

\subsection{Test error, allocations and arithmetic}
\label{app:comparison-opt}

For a test probability measure $Q$, define
\begin{equation}
 E_Q(g;T)=\sup_{t\le T}\norm{g(t)-f_\infty(t)}_{L^2(Q)},\qquad
 \mathcal E_Q(g;T)=\bigl(\E E_Q(g;T)^2\bigr)^{1/2}.
 \label{cmp:test-error}
\end{equation}
The reverse triangle inequality bounds the difference between the test RMSEs
of $g$ and $f_\infty$, against any $y_*\in L^2(Q)$, by $E_Q(g;T)$.
This measures approximation, without assuming that $f_\infty$ generalizes well.
Moreover, pointwise estimates
$\E\sup_{t\le T}|f_n(t,x)-f_\infty(t,x)|^2\le C_T(x)^2/n$
with $C_T\in L^2(Q)$ give
$\mathcal E_Q(f_n;T)^2\le n^{-1}\int C_T^2\dd Q$
by moving the time supremum inside the integral and applying Tonelli.
An input supremum would require additional spatial control.

For $\phi=\tanh$, set
$a_\ell=\norm{h^{(\ell)}-h'^{(\ell)}}_2/\sqrt n$ and assume
$\norm{x}/\sqrt d\le X$.  Boundedness and Lipschitz continuity of $\tanh$ give
\begin{align*}
 a_1&\le X\frac{\norm{W^{(1)}-W'^{(1)}}_F}{\sqrt n},\qquad
 a_\ell\le\norm{W^{(\ell)}-W'^{(\ell)}}_{\rm op}
             +\norm{W'^{(\ell)}}_{\rm op}a_{\ell-1},\\
 |f_n-f'_n|&\le\frac{\norm{w-w'}_2}{\sqrt n}
                    +\frac{\norm{w'}_2}{\sqrt n}a_H.
\end{align*}
Induction transfers parameter control to every such test input when operator
and normalized readout norms are bounded.  Coupling dense and compressed
networks through their initialization, the triangle inequality and Minkowski give
\begin{equation}
 \mathcal E_Q(\widehat f_{n,P};T)
 \le\mathcal E_Q(f_n;T)
 +\left(\E\sup_{t\le T}\norm{\widehat f_{n,P}(t)-f_n(t)}_{L^2(Q)}^2\right)^{1/2}.
 \label{cmp:triangle}
\end{equation}
Thus the conditional model $A n^{-1/2}+B_qP^{-q}$, $q=1,2$, needs both a
matched root-width estimate and closure constants and sufficient-order
thresholds uniform in width, with appropriate moments.  The perturbative
fluctuation calculation of \citet{bordelon2023finite} motivates the first
estimate but is not a general remainder theorem for this model.

Moving memory counts persistent updated scalars; fixed matrices, coefficient
evaluators and workspace are reported separately.  The leading response-memory
count is $2(H-1)mnP$.  With $u=A/\sqrt n$, $v=B_q/P^q$ and $u+v=\varepsilon$,
its minimization reduces to minimizing $u^{-2}(\varepsilon-u)^{-1/q}$.
The logarithm has derivative $-2/u+1/[q(\varepsilon-u)]$ and positive second
derivative, and diverges at both endpoints.  Hence the unique continuous optimum is
\begin{equation}
 u_*=\frac{2q\varepsilon}{2q+1},\quad v_*=\frac{\varepsilon}{2q+1},\quad
 n_*=\frac{(2q+1)^2A^2}{4q^2\varepsilon^2},\qquad
 P_* =\left(\frac{(2q+1)B_q}{\varepsilon}\right)^{1/q}.
 \label{cmp:allocation}
\end{equation}
For $q=1$, $(n_*,P_*)=(9A^2/(4\varepsilon^2),3B_1/\varepsilon)$ and
$M_*\sim\tfrac{27}{2}(H-1)mA^2B_1\varepsilon^{-3}$.
For $q=2$, $(n_*,P_*)=(25A^2/(16\varepsilon^2),\sqrt{5B_2/\varepsilon})$ and
$M_*\sim\tfrac{25\sqrt5}{8}(H-1)mA^2\sqrt{B_2}\varepsilon^{-5/2}$.
Upward rounding is feasible.  Outer weights, $(d+1)n$, and the joint Gram
matrix, $O(P^2)$, perturb the finite-tolerance optimum but are lower order.
Dense training instead takes $n=\lceil(A/\varepsilon)^2\rceil$ and
$M_*\sim(H-1)A^4\varepsilon^{-4}$.  These optimize sufficient bounds in
specified representations, not actual errors or all possible algorithms.

For completeness, full-batch arithmetic per right-hand-side evaluation is
\begin{align}
 C_{\rm dense}&=O(Hmn^2+mnd),\nonumber\\
 C_{\rm old}&=O(Hmn^2+Hm^2nP+HmnP+mnd),\nonumber\\
 C_{\rm joint}&=O(Hmn^2+Hm^2nP+HmnP^2+P^3+mnd).
 \label{cmp:response-work}
\end{align}
The terms respectively account for dense initialized actions, learned factors,
moment updates or Gram solves, and Gram factorization.  Directional passes for
the joint clock add a constant number of network passes.  Both closures retain
$(H-1)n^2$ fixed entries; activation workspace is $O(Hmn)$, and Gram
conditioning can affect numerical step counts.  Consequently neither their
total-storage nor their arithmetic exponent improves automatically.

\subsection{A matched tangent hierarchy}
\label{app:comparison-nth}

With $\mathcal D_n=\diag(nI,I,\ldots,I,nI)$, define recursively
\begin{align*}
K^{(2)}(x,z)&=\nabla_\theta f(x)^\top\mathcal D_n\nabla_\theta f(z),\\
K^{(j+1)}(x_1,\ldots,x_{j+1})&=
(\nabla_\theta K^{(j)}(x_1,\ldots,x_j))^\top
\mathcal D_n\nabla_\theta f(x_{j+1}).
\end{align*}
The chain rule gives the exact, mobility-matched hierarchy
\begin{equation}
 \dot f(x)=-\frac{2}{m}\sum_a r_aK^{(2)}(x,x_a),\qquad
 \dot K^{(j)}=-\frac{2}{m}\sum_a r_aK^{(j+1)}(\ldots,x_a).
 \label{cmp:nth-flow}
\end{equation}
Truncation through order $p\ge2$ retains exact initial coefficients and freezes
$K^{(p)}$.  A population version additionally requires limiting coefficients
and a valid limiting hierarchy.

\paragraph{Feedback bound.}
Suppose $|K^{(j)}|\le B_j$ along the exact flow, uniformly when the first
argument is a training or test input and the others are training inputs.
Let $R=1+\sup_{t,a}|r_a(t)|$ and stop the truncated system at predictor error
one.  Its residuals are then bounded by $R$.  For supremum predictor error
$e_f$ and order-$j$ kernel error $e_j$, subtraction, splitting each product
using the exact kernel and truncated residual, yields
\begin{align}
 e_p(t)&\le2RB_{p+1}t,\nonumber\\
 e_j(t)&\le\int_0^t(2Re_{j+1}+2B_{j+1}e_f)\dd s\quad(2\le j<p),\nonumber\\
 e_f(t)&\le\int_0^t(2Re_2+2B_2e_f)\dd s.
 \label{cmp:nth-recursion}
\end{align}
Successive substitution integrates the top forcing $p$ times.  Each feedback
term has convolution kernel $2B_{k+2}(2R(t-s))^k/k!$, $0\le k\le p-2$.
Bounding these kernels at $T$ and applying Gronwall gives
\begin{equation}
 \sup_{t\le T}e_f(t)\le\eta_p(T):=
 B_{p+1}\frac{(2RT)^p}{p!}\exp(L_{p,T}T),\qquad
 L_{p,T}=2\sum_{k=0}^{p-2}B_{k+2}\frac{(2RT)^k}{k!}.
 \label{cmp:nth-bound}
\end{equation}
If $\eta_p(T)<1$, first exit is impossible.  The triangular integral equations
then bound all training tensors, preventing finite-time escape and extending
the solution through $T$.  The test bound implies $E_Q\le\eta_p$ for $Q$
supported on the controlled test set.  For random coefficients, the same
probabilistic convention and coefficient errors must also be controlled.

One sufficient regime is $B_j\le UV^{j-2}(j-2)!$ with $U,V>0$ and
$\chi=2RVT<1$.  Then $L_{p,T}\le2U/(1-\chi)$ and
$\eta_p\le D\chi^p/p$, where $D=(U/V)e^{2UT/(1-\chi)}$.
For $0<\varepsilon<1$, put $\lambda_\chi=\log(1/\chi)$.  The smallest
certified integer order is
\begin{equation}
 p_\varepsilon=\max\left\{2,\left\lceil
 \frac{W(\lambda_\chi D/\varepsilon)}{\lambda_\chi}\right\rceil\right\},\qquad
 M_{\rm NTH}=\sum_{j=1}^{p_\varepsilon-1}m^j
 =\frac{m^{p_\varepsilon}-m}{m-1}\quad(m\ge2).
 \label{cmp:nth-memory}
\end{equation}
Here $W(z)e^{W(z)}=z$: the error inequality is equivalent to
$(\lambda_\chi p)e^{\lambda_\chi p}\ge\lambda_\chi D/\varepsilon$.
The top tensor has $m^{p_\varepsilon}$ \emph{fixed} entries, and direct work
is $O(m^{p_\varepsilon})$ per evaluation.  Depth is already combined into
these kernels.  For fixed $m\ge2$ and $\chi$, the moving count is
$\asymp[D/(\varepsilon\log(D/\varepsilon))]^{\log m/\lambda_\chi}$.
Tanh analyticity alone does not supply these uniform kernel bounds;
$\chi\ge1$ makes this certificate inconclusive.  Population use requires
$U,V,R$ independent of the approximation resources.  Finite-width initial
coefficients also require a width-error allowance; with population
coefficients, their construction and approximation are fixed-resource costs.

\paragraph{Scope of the published estimate.}
After exchanging their width/sample symbols, \citet[Theorem~2.6]{huang2020nth}
give the training RMS bound
\begin{align}
 \frac{\norm{f(t)-\widetilde f(t)}_2}{\sqrt m}
 &\lesssim\frac{(1+t)t^{p-1}}{n^{p/2}}\min\{t,m/\lambda\},\nonumber\\
 t&\le\min\left\{\frac{c\sqrt{\lambda n/m}}{(\log n)^C},
 \frac{n^{p/[2(p+1)]}}{(\log n)^{C'}}\right\}.
\label{cmp:published-nth}
\end{align}
Time uses the source loss convention; a factor-two time rescaling changes
constants, not these powers.
This uses their high-probability initialization setting and parameterization,
even $p$, bounded
activation derivatives through $2p+1$, nondegenerate bounded inputs with
uniformly independent subsets through size $2p+1$, and an initial kernel
eigenvalue exceeding $\lambda$.  Three vectors in two dimensions violate the
subset assumption.  Their small higher-kernel estimates do not establish
\eqref{cmp:nth-bound}'s population hypotheses for our feature-learning flow;
their training RMS estimate is not a whole-test-function estimate.

\paragraph{Whole-function evaluation and time analyticity.}
Set $u_a=-2\widetilde r_a/m$, $I_\varnothing=1$, and
$\dot I_{a_1\cdots a_k}=u_{a_1}I_{a_2\cdots a_k}$ with zero initial values
for nonempty words.  Repeated integration of the triangular system gives
\begin{equation}
 \widetilde f_p(t,x)=f_0(x)+\sum_{k=1}^{p-1}\sum_{a_1,\ldots,a_k}
 K_0^{(k+1)}(x,x_{a_1},\ldots,x_{a_k})I_{a_1\cdots a_k}(t).
 \label{cmp:signature}
\end{equation}
This uses the same number of moving coefficients as the direct hierarchy;
one query needs that many contractions \emph{plus} the cost of the fixed
mixed-kernel evaluators.  When $m=1$, $I_k=s^k/k!$ and $\dot s=u$, so one
moving scalar suffices and order dependence remains in fixed coefficients
and evaluation.  These counts are therefore representation costs, not lower bounds.

Truncation is not a time-Taylor polynomial.  For the scalar model
$f(\theta)=\theta^2$, loss $f^2$, target zero and unit mobility,
$f_0>0$ gives $\dot\theta=-4\theta^3$, $f(t)=f_0/(1+8f_0t)$ and
$K^{(j)}=4^{j-1}f$.  Its time-Taylor radius is $1/(8f_0)$, but
\eqref{cmp:signature} becomes
$\widetilde f_p=f_0\sum_{k=0}^{p-1}(4s)^k/k!$, $\dot s=-2\widetilde f_p$.
The exact scalar field replaces the sum by $e^{4s}$, with solution
$s=-\tfrac14\log(1+8f_0t)$.  On a compact neighborhood of this path for any
finite $T$, the polynomial fields and their derivatives converge uniformly.
Subtracting the ODEs gives a vanishing forcing and a uniform Lipschitz
coefficient; Gronwall and first exit give convergence on $[0,T]$.
Thus a finite time-Taylor radius is no universal barrier for NTH.

\subsection{A specified streamed DMFT solver}
\label{app:comparison-dmft}

The sampled self-consistency algorithm of
\citet[Appendix~B, Algorithm~1]{bordelon2022dmft} stores layerwise covariance
and response kernels indexed by input and time pairs.  Each sweep updates
predictions, factors Gaussian covariances, samples paths, solves their causal
equations and response Jacobians, and averages products and derivatives.
Matching its feature-learning scale to ours requires $\gamma=\sqrt n$,
internal coordinates $U^{(\ell)}=\sqrt nW^{(\ell)}$, and the same readout
law and loss normalization.  A fixed-program Tensor Programs limit
\citep[Algorithm~1, Theorem~7.4]{yang2021tensor} does not itself prescribe
this numerical solver or a uniform numerical error estimate.

For $K$ time steps, $S$ paths and $I$ sweeps, a dense covariance factor costs
$O((mK)^3)$ per layer.  Dense triangular solves for a full samplewise
response Jacobian cost at most the same per path.  Streaming one path and
Jacobian while accumulating kernel averages therefore permits
\begin{equation}
 M_{\rm D}=O(Hm^2K^2),\qquad
 C_{\rm D,total}=O(IH(S+1)m^3K^3),\qquad
 C_{\rm D,step}=O(IH(S+1)m^3K^2).
 \label{cmp:dmft-cost}
\end{equation}
The last cost is amortized; memory includes kernel and Jacobian workspace.
These are implementable upper bounds, not framework-wide lower bounds.
Passive test prediction requires mixed kernels and test paths as well.

To justify an error model, suppose the discrete self-consistency map
$\mathcal F$ is a $\kappa$-contraction, $\kappa<1$, on a common domain,
in a norm controlling test outputs, has fixed point $z_*$, and the sampled map
$\widehat{\mathcal F}$ differs
uniformly by at most $\zeta$.  For a computed point $\widehat z$ with residual
$\norm{\widehat z-\widehat{\mathcal F}(\widehat z)}\le\eta_0$, the triangle
inequality gives
$\norm{\widehat z-z_*}\le\eta_0+\zeta+\kappa\norm{\widehat z-z_*}$.
A Lipschitz output map, uniform RMS sampling control and a stable order-$r$
time approximation would thus yield
\begin{equation}
 \mathcal E_Q\le A_{\rm D}S^{-1/2}+B_{\rm D}(T/K)^r+\eta,
 \qquad \eta\le C\eta_0/(1-\kappa).
 \label{cmp:dmft-error}
\end{equation}
The constants include stability amplification.  The cited papers do not
establish these hypotheses uniformly as the grid grows for this model.
Setting $S=n$ matches a sampling exponent, not the covariance of finite-network
fluctuations \citep{bordelon2023finite}.

Since streaming makes leading memory independent of $S$, its continuous
memory infimum has $K_{\inf}=T(B_{\rm D}/\varepsilon)^{1/r}$ and
$M_{\inf}\asymp Hm^2T^2B_{\rm D}^{2/r}\varepsilon^{-2/r}$.
A finite feasible allocation reserves $\beta\varepsilon$, $0<\beta<1$,
for sampling and iteration: take
$K=\lceil T[B_{\rm D}/((1-\beta)\varepsilon)]^{1/r}\rceil$,
$S\ge(2A_{\rm D}/(\beta\varepsilon))^2$, and
$\eta\le\beta\varepsilon/2$.  Approaching the infimum increases work.
Optimizing numerical order requires additional regularity, constants and
stage costs; there is no intrinsic DMFT accuracy exponent here.

A distinct finite-width baseline stores every Euler update exactly:
\[
 W^{(\ell)}_k=W^{(\ell)}_0-\frac{2\Delta}{nm}
 \sum_{j<k,a}r_{a,j}\delta^{(\ell)}_{a,j}(h^{(\ell-1)}_{a,j})^\top.
\]
Its moving memory is $O(HmnK)$, fixed memory $O(Hn^2)$ and late-step work
$O(Hmn^2+Hm^2nK)$.  A fixed-stage order-$r$ method has the same storage
order; with error $A/\sqrt n+D_{\rm h}(T/K)^r$, optimizing as above gives
$M_*\asymp HmA^2T D_{\rm h}^{1/r}\varepsilon^{-(2+1/r)}$.
This stores finite-network updates and is not Tensor Programs sampling.

\subsection{A sharper estimate for the original joint clock}
\label{app:comparison-joint}

Keep exactly $\dot\tau=\rho+\norm{\dot\Psi}_2$, with its matching prefix.
Define
\[
M_T=1+\int_0^T\rho\dd t,\qquad
V_T=\frac{1}{\sqrt{nm}}\int_0^T\norm{\dot\Psi}_2\dd t,\qquad
\tau_T=M_T+\sqrt{nm}V_T.
\]
The insertion measure has mass $M_T$, while the normalized history
$F(\xi)=\Psi(\xi)/\sqrt{nm}$ is $1/\sqrt{nm}$-Lipschitz.

We need a dimension-independent polynomial approximation fact.  If a
Hilbert-valued $F$ is $L$-Lipschitz on $[0,\tau]$, then some polynomial
of degree below $P$ has uniform error at most $C_JL\tau/P$.
To prove it, extend $G(v)=F(\tau(1+\cos v)/2)$ evenly and periodically and
convolve with the normalized positive trigonometric polynomial
$J_N(v)=Z_N^{-1}[\sin(Nv/2)/\sin(v/2)]^4$.
On $[-\pi,\pi]$ the ratio is bounded by $\min(N,\pi/|v|)$ and is at least
$cN$ near zero.  Splitting its integrals at $1/N$ gives $Z_N\asymp N^3$
and $\int|v|J_N(v)\dd v\le C/N$.
The Hilbert norm triangle inequality bounds convolution error by
$C L\tau/N$.  The convolution is even, hence a polynomial in $\cos v$
of degree at most $2N-2$.  Take $N=\lfloor(P+1)/2\rfloor$ for $P\ge2$;
a constant polynomial handles $P=1$.

The weighted projection is no worse than this uniform approximation in
$L^2(\mu_T)$.  Using the projection energies and product defect of
\eqref{eq:defect-energy} therefore gives
\begin{equation}
 \int_0^T\sum_{\ell=2}^H\norm{E_\ell}_F\dd t
 \le \frac{1}{nm}\sum_{\ell,a}(D_{h,\ell,a}+D_{b,\ell,a})
 \le \frac{C_J^2M_T}{P^2}
       \left(\frac{M_T}{\sqrt{nm}}+V_T\right)^2.
 \label{cmp:weighted-jackson}
\end{equation}
The first inequality is Cauchy--Schwarz in time followed by
$2\sqrt{ab}\le a+b$; the second uses the actual insertion mass rather than
Lebesgue mass $\tau_T$.  No clock or algorithm has changed.  Width-uniform
history mass, normalized variation, feedback stability and continuation
thresholds would turn this into a uniform prediction estimate.  Establishing
those controls, and the matched width rate, remains the population gap.
